\documentclass[12pt, a4paper]{article}

% --- PREAMBOLO ---
\usepackage[utf8]{inputenc}
\usepackage[italian]{babel}
\usepackage[T1]{fontenc}
\usepackage{amsmath, amssymb}
\usepackage{graphicx}
\usepackage[utf8]{inputenc}
\usepackage[T1]{fontenc}
\usepackage{media9}
\usepackage{graphicx}
\usepackage{geometry}
\geometry{margin=2.5cm}

% Comando per la linea orizzontale nella copertina
\newcommand{\hr}{\rule{\linewidth}{0.5mm}}

% --- INIZIO DOCUMENTO ---
\begin{document}
\begin{titlepage}
    \centering
    
    \begin{minipage}{0.20\textwidth}
        \centering
        
\includegraphics[width=\textwidth]{immagine dii}
    \end{minipage}
    \hfill
    \begin{minipage}{0.55\textwidth}
        \centering
        {\Large \textbf{Università degli Studi di Padova}} \\
        \vspace{0.3cm}
        {\large Corso di Laurea in Ingegneria Aerospaziale} \\
        \vspace{0.3cm}
        {\normalsize Canale: B}
    \end{minipage}
    \hfill
    \begin{minipage}{0.20\textwidth}
        \centering
        
\includegraphics[width=\textwidth]{Logo Università Padova} 
    \end{minipage}

    \vspace{1cm}
    
    \hrule 
    \vspace{0.6cm}
    {\Huge \bfseries Progettino Dinamica del Volo} \\
    \vspace{0.4cm}
    {\large Prima Parte: Descrizione Orbita}
    \vspace{0.6cm}
    \hrule
    
    \vspace{2cm}
 
    \includegraphics[width=0.8\textwidth]{"orbita"} 
    
    \vfill 
   
    
    \vfill 
    
    \begin{minipage}{0.5\textwidth}
        \begin{flushleft} \large
            \emph{Autori:}\\
            Pietro \textbf{Guglielmini} (Matr. 2140457)\\
            Luigi \textbf{Ippolito} (Matr. 2143068)\\
            Alessandro \textbf{Magentini} (Matr. 2139775)\\
            Filippo \textbf{Michielan} (Matr. 2138899)\\
            
            
        \end{flushleft}
    \end{minipage}
    \begin{minipage}{0.4\textwidth}
        \begin{flushright} \large
            \emph{Docente:} \\
            Prof. Giacomo \textbf{Colombatti}
        \end{flushright}
    \end{minipage}
    
    \vspace{1.5cm}
    {\large \today}
\end{titlepage}

\newpage
\tableofcontents 
\newpage

% --- CONTENUTO ---

\section{Introduzione}
Il primo progettino relativo al corso di Dinamica del Volo, anno accademico 2025/2026, richiede ad un gruppo di massimo quattro studenti di disegnare una missione di intercettazione di un asteroide posizionato su di un' orbita con inclinazione relativa rispetto al piano equatoriale terrestre.
\section{Scelta Software specifico}
Il software specifico è stato scelto attraverso interrogazione dell'AI Claude e Gemini. Il prompt utilizzato è il seguente, riportato in forma integrale:  

\begin{quote}
    \itshape % Rende il font corsivo
    \textsf{ % Rende il font Sans Serif r
        ``Devo svolgere un progetto universitario per il corso di Dinamica del Volo Aerospaziale in un gruppo di 4 persone. L'obiettivo è disegnare una missione di intercettazione di un asteroide su orbita inclinata rispetto al piano equatoriale terrestre. I contenuti richiesti includono: parametri orbitali, posizione del satellite ogni 3 ore, plot 3D in sistema ECI, analisi delle perturbazioni, ground track, manovre orbitali e phasing. La metodologia richiede analisi teorica, motivazione degli strumenti usati, rilevazione delle discrepanze tra teoria e simulazione, e identificazione del contributo di ogni membro. Il professore propone come alternative MATLAB o GMAT. Analizza pro e contro di entrambi rispetto a questo specifico progetto, valuta se convenga usarli entrambi in parallelo o in modo complementare. ''
    }
\end{quote}

La risposta fornita è stata la seguende :

\begin{quote}
    \itshape  % Attiva il corsivo per lo stile citazione
    \sffamily % Attiva il font Sans Serif in modo corretto per testi lunghi
    
    \textbf{MATLAB} \\
    Pro: conosciuto da tutti, lavoro divisibile per moduli (propagatore, perturbazioni, manovre, plot), codice commentato significa documentazione trasparente, facile mostrare contributo individuale e corrispondenza teoria, implementazione. \\
    Contro: perturbazioni (J2, drag, SRP) e manovre da implementare a mano, validazione dei risultati interamente a carico nostro.

    \medskip 

    \textbf{GMAT} \\
    Pro: strumento NASA nativo per questo tipo di missione, perturbazioni, ground track, manovre, phasing già integrati e validati, dà peso professionale al report. \\
    Contro: curva di apprendimento non banale (interfaccia, scripting proprio, debug difficile), corrispondenza teoria, implementazione difficile da documentare, meno flessibile per visualizzazioni personalizzate.

    \medskip

    \textbf{Entrambi in parallelo (2+2 persone)} \\
    Pro: confronto diretto tra i due approcci, ognuno lavora con lo strumento preferito. \\
    Contro: overhead di coordinamento alto per convergere su un unico report coerente, sforzo molto non bilanciato tra i due sottogruppi (chi usa MATLAB avanza molto più velocemente), nessun vantaggio metodologico esplicito richiesto dal prof.

    \medskip

    \textbf{MATLAB principale + GMAT per validazione} \\
    Pro: carico limitato su GMAT (solo run di confronto su propagazione e manovre chiave), risultati validati da riferimento esterno, diventa naturalmente la sezione sulle discrepanze teoria, simulazione richiesta dalla metodologia. \\
    Contro: richiede che almeno uno impari GMAT a sufficienza per configurare e interpretare i run, tempi da coordinare con il resto del gruppo.
\end{quote}
Analizzata la situazione abbiamo scelto di utilizzare MATLAB come software principale per l'analisi dell'orbita mentre GMAT è stato utilizzato come confronto e validazione dei dati ottenuti. 

\begin{figure}[htbp]
    \centering
    \includegraphics[width=0.8\textwidth]{orbita.png}
    \caption{Rappresentazione dell'orbita di missione} 
    \label{fig:mio_grafico}
\end{figure}
\section{Analisi del problema e Implementazione codice}
Il primo passaggio dello studio è stato analizzare l'orbita attraverso GMAT in modo tale da avere una base sulla quale verificare i dati ricavati successivamente. Una volta ottenuto questo primo riscontro siamo passati ad implementare il codice MATLAB, codice in grado di descrivere e rappresentare l'orbita richiesta. Per analizzare i risultati sono stati utilizzati i toll specifici Matlab "aerospace toolbox" e "satellite comunication toolbox".
All'interno del codice MATLAB sono state implementate diverse funzioni: una per il calcolo dei parametri orbitali e una per la propagazione. Quest'ultima, a partire dall'intervallo temporale definito inizialmente, segue le equazioni descritte in Sezione \ref{sec:formule} e calcola l'orbita seguendo l'ordine t - $M_e$ - xi - TA - r. Questa metodologia però, nonostante fornisca risultati corretti, non permette una visualizzazione dell'orbita nel sistema ECI, ma soltanto nel sistema NEU (North East Up) in MATLAB lla (Latitude Longitude Altitude), dunque la visualizzazione dell'orbita è possibile soltanto da un punto di vista fisso con la rotazione terrestre. Per ottenere l'orbita nel sistema ECI, quindi, ci siamo appoggiati alle functions del Satellite Comunication Toolbox affiancato al Aerospace Toolbox, in particolare \texttt{satellite}, che propaga l'orbita, e \texttt{satelliteScenarioViewer} che visualizza l'orbita nella maniera da noi desiderata.
Nello specifico il modello implementato dall'Aerospace Toolbox esegue un modello simulink. 

\section{Algoritmo di Determinazione dei Parametri Orbitali}
\label{sec:formule}
Partendo dai vettori stato $\vec{r}$ e $\vec{v}$, l'algoritmo calcola i parametri orbitali classici (COE). Si definisce il versore dell'asse $z$ come $\hat{k} = [0, 0, 1]^T$.

\subsection{Costanti del Moto}
Il vettore momento angolare specifico $\vec{h}$ e l'energia specifica $\varepsilon$ sono calcolati come:
\begin{align}
    \vec{h} &= \vec{r} \times \vec{v} \\
    \varepsilon &= \frac{v^2}{2} - \frac{\mu}{r}
\end{align}

\subsection{Parametri per Orbita Ellittica ($\varepsilon < 0$)}
Se l'energia specifica è negativa, l'orbita è di tipo ellittico. I parametri vengono determinati come segue:

\begin{itemize}
    \item \textbf{Semiasse maggiore e Periodo:}
    \begin{equation}
        a = -\frac{\mu}{2\varepsilon}, \quad T = \frac{2\pi}{\sqrt{\mu}} a^{3/2}
    \end{equation}

    \item \textbf{Vettore Eccentricità:}
    \begin{equation}
        \vec{e} = \frac{1}{\mu} \left[ \vec{v} \times \vec{h} \right] - \frac{\vec{r}}{r}, \quad e = \|\vec{e}\|
    \end{equation}

    \item \textbf{Inclinazione ($i$):}
    L'inclinazione è l'angolo tra il momento angolare $\vec{h}$ e l'asse $z$:
    \begin{equation}
        i = \arccos\left( \frac{h_z}{h} \right)
    \end{equation}

    \item \textbf{Ascensione Retta del Nodo Ascendente ($\Omega$ - RAAN):}
    Si definisce il vettore linea dei nodi $\vec{N} = \hat{k} \times \vec{h}$.
    \begin{equation}
        \Omega = 
        \begin{cases} 
        \arccos\left(\frac{N_x}{N}\right) & \text{se } N_y > 0 \\
        2\pi - \arccos\left(\frac{N_x}{N}\right) & \text{se } N_y \leq 0
        \end{cases}
    \end{equation}

    \item \textbf{Argomento del Pericentro ($\omega$):}
    L'angolo tra la linea dei nodi $\vec{N}$ e il vettore eccentricità $\vec{e}$:
    \begin{equation}
        \omega = 
        \begin{cases} 
        \arccos\left(\frac{\vec{N} \cdot \vec{e}}{N e}\right) & \text{se } e_z > 0 \\
        2\pi - \arccos\left(\frac{\vec{N} \cdot \vec{e}}{N e}\right) & \text{se } e_z \leq 0
        \end{cases}
    \end{equation}
\end{itemize}

\section{Propagazione Temporale e Moto di Keplero}

Per determinare l'evoluzione temporale della posizione sulla traiettoria, l'algoritmo risolve il problema del tempo di volo partendo dall'anomalia media fino alla determinazione delle coordinate nello spazio ECI (\textit{Earth-Centered Inertial}).

\subsection{Definizione delle Variabili Temporali}
Dato il periodo orbitale $T$ e un passo temporale di campionamento $dt$, si definisce il numero di intervalli $n = \lfloor T/dt \rfloor + 1$. L'anomalia media $M_e$ per ogni istante $t_k$ è calcolata come:
\begin{equation}
    M_e(k) = \frac{2\pi}{T} \cdot t_k, \quad \text{con } t_k = (k-1)dt
\end{equation}

\subsection{Risoluzione dell'Equazione di Keplero}
Il passaggio dall'anomalia media $M_e$ all'anomalia eccentrica $E$ (indicata nel codice come $\xi$) richiede la risoluzione della seguente equazione trascendente:
\begin{equation}
    f(E) = E - e \sin(E) - M_e = 0
\end{equation}
L'algoritmo utilizza un risolutore numerico (\textit{fzero}) con una stima iniziale $E_0$ ottimizzata per migliorare la convergenza:
\begin{equation}
    E_0 = 
    \begin{cases} 
    M_e + e/2 & \text{se } M_e \leq \pi \\
    M_e - e/2 & \text{se } M_e > \pi
    \end{cases}
\end{equation}

\subsection{Determinazione della Posizione Orbitale}
Una volta ottenuta l'anomalia eccentrica $E$, si procede al calcolo dell'anomalia vera $\theta$ (TA nel codice) e del raggio orbitale $r$:
\begin{align}
    \theta &= 2 \arctan \left( \sqrt{\frac{1+e}{1-e}} \tan\frac{E}{2} \right) \\
    r(\theta) &= \frac{h^2/\mu}{1 + e \cos \theta}
\end{align}

\subsection{Trasformazione di Coordinate (da Kepleriano a ECI)}
L'ultimo passaggio dell'algoritmo consiste nella rotazione dal piano orbitale alla terna inerziale $IJK$ (ECI). Utilizzando i parametri orbitali $i$ (inclinazione), $\Omega$ (RAAN) e $\omega$ (argomento del pericentro), la posizione $\vec{r}_{ECI}$ e la velocità $\vec{v}_{ECI}$ sono ottenute tramite la matrice di rotazione $\mathbf{R}(\Omega, i, \omega)$:
\begin{align}
    \vec{r}_{ECI} &= \mathbf{R}_{z}(\Omega) \mathbf{R}_{x}(i) \mathbf{R}_{z}(\omega) \vec{r}_{perifocale} \\
    \vec{v}_{ECI} &= \mathbf{R}_{z}(\Omega) \mathbf{R}_{x}(i) \mathbf{R}_{z}(\omega) \vec{v}_{perifocale}
\end{align}
Dove le componenti nel sistema perifocale dipendono esclusivamente da $r, \theta, \mu, p$.

\subsection{Note sull'implementazione}
Tutti i valori angolari calcolati in radianti sono stati successivamente convertiti in gradi per facilitare l'analisi dei risultati:
\begin{equation}
    \text{valore [deg]} = \text{valore [rad]} \cdot \frac{180}{\pi}
\end{equation}

\section{Analisi dei Risultati}

\section{Organizzazione del gruppo}
Il lavoro del gruppo è stato organizzato attorno alla piattaforma GitHub, scelta che abbiamo preferito rispetto ad una più tradizionale e diffusa cartella condivisa su un drive online, in quanto permette di gestire meglio e di tenere traccia di tutte le modifiche ai codici, inoltre non presenta limiti in quanto a spazio di archiviazione. \'E stata creata una repository alla quale tutti i componenti del gruppo hanno avuto accesso permettendo di collaborare, e si è optato per lavorare tutti su una sola branch visto che complicare la struttura della repository non avrebbe portato ad alcun vantaggio aggiuntivo. 
I file sono stati gestititi in due cartelle separate: per i codici (sia GMAT che MATLAB) è stata creata l'apposita cartella \textsl{Codici}, mentre il codice LaTeX e le immagini sono stati inseriti nella cartella \textsl{Documentazione}.
In questa prima fase del progetto i componenti del gruppo Alessandro Magentini e Filippo Michielan si sono dedicati alla parte di programmazione software, in particolare all'implementazione prima GMAT e poi MATLAB dell'orbita, mentre Ippolito Luigi e Guglielmini Pietro hanno svolto l'analisi dati e la stesura del report in LaTeX. Il gruppo ha lavorato in maniera sincrona  con incontri in presenza e su piattaforme online come google meet e discord, ma anche asincrona sui programmi specifici. 

\end{document}